\section{具有其他单规范群的理论}\label{sec:higherrank}
我们已经花了相当多的篇幅研究规范群为 $\SU(2)$ 的 $\cN{=}2$ 规范理论。本节转向对更大规范群的分析。首先较为详细地研究 $\SU(N)$ 规范理论，然后讨论 $\SO(2N)$ 的情形。我们还将分析由这些规范理论得到的 Argyres-Douglas 共形场论，并说明它们由第 \ref{sec:generalAD} 小节引入的理论 $X_N$ 和 $Y_N$ 给出。最后，在第 \ref{sec:others} 小节简要提及具有其他规范群的理论的 Seiberg-Witten 解，结束本节。

本节将给出的曲线，可能并不是较早文献中最常见的形式。它们之间的关系也将在第 \ref{sec:others} 小节说明。

\subsection{半经典分析}
考虑含有 $N_f$ 个基本 $N$ 维表示超多重态的 $\SU(N)$ 规范理论。$\cN{=}2$ 矢量多重态由 $\cN{=}1$ 伴随表示手征多重态 $\Phi$ 和 $\cN{=}1$ 矢量多重态 $W_\alpha$ 组成，两者都是 $N\times N$ 矩阵。用 $\cN{=}1$ 手征多重态表示时，超多重态可以写为
\begin{equation}
Q_i^a,\quad\tilde Q^i_a,\qquad a=1,\ldots,N;\quad i=1,\ldots,N_f.
\end{equation}
超对称真空的一个分支由条件
\begin{equation}
[\Phi,\Phi^\dagger]=0.
\end{equation}
给出。这意味着 $\Phi$ 可以对角化。将其记为
\begin{equation}
\Phi=\diag(a_1,\ldots,a_N),\qquad\sum a_i=0.
\end{equation}

考虑一般的情形，即对所有 $i\neq j$ 都有 $a_i\neq a_j$。此时规范群从 $\SU(N)$ 破缺到 $\U(1)^{N-1}$。来自 $N\times N$ 矩阵的 $(i,j)$ 元的 W 玻色子，其质量为
\begin{equation}
M_W=|a_i-a_j|.
\end{equation}
至于单极子，已知从 $\SU(2)$ 破缺到 $\U(1)$ 的 't Hooft-Polyakov 单极子解，可以直接视为从 $\SU(N)$ 破缺到 $\U(1)$ 的解：只需在 $N\times N$ 矩阵中选取 $2\times2$ 子矩阵，即对于 $i\neq j$，取出位于 $(i,i)$、$(i,j)$、$(j,i)$ 和 $(j,j)$ 的矩阵元。此时质量为
\begin{equation}
M_\text{monopole}=|\tau_{UV}(a_i-a_j)|.
\end{equation}

单圈跑动为
\begin{equation}
\LambdaRG\frac d{d\LambdaRG}\tau=\frac1{2\pi}(2N-N_f).
\end{equation}
因此，当 $0\le N_f<2N$ 时，理论是渐近自由的。此时动力学能标为
\begin{equation}
\Lambda^{2N-N_f}:=\Lambda_{UV}^{2N-N_f}e^{2\pi i\tau_{UV}}.
\end{equation}
当 $N_f=2N$ 时，理论是渐近共形的，$\tau_{UV}$ 是量子理论中的一个无量纲参数。

存在味时，这一真空中的 $\cN{=}1$ 超势为
\begin{multline}
\sum_i(Q_i\Phi\tilde Q^i-\mu_iQ_i\tilde Q^i)
=\\
\sum_i(Q_i^1,Q_i^2,\ldots,Q_i^N)
\begin{pmatrix}
a_1-\mu_i\\
&a_2-\mu_i\\
&&\ddots\\
&&&a_N-\mu_i
\end{pmatrix}
\begin{pmatrix}
\tilde Q^i_1\\
\tilde Q^i_2\\
\vdots\\
\tilde Q^i_N
\end{pmatrix}.
\end{multline}
于是，只要对某些 $i$ 和 $s$ 有 $a_i-\mu_s=0$，就会出现一个无质量的带荷超多重态分量。

在强耦合量子理论中，将 $a_i$ 定义为依赖规范的量 $\Phi$ 的对角元，并没有多大意义。相反，正如在 $\SU(2)$ 规范理论中所做的那样，将 $a_i$ 定义为进入 BPS 质量公式的复数：
\begin{equation}
M\ge|n^ia_i+m_ia_D^i+\sum_sf_s\mu_s|
\end{equation}
其中 $(n^i,m_i)$ 是红外规范群 $\U(1)^{N-1}$ 下的电荷与磁荷，$f_s$ 是味荷。还可以考虑由下式定义的 $\Phi$ 的规范不变组合：
\begin{equation}
x^N+u_2x^{N-2}+\cdots+u_{N-1}x+u_N:=\vev{\det(x-\Phi)}
\end{equation}
其中 $x$ 是一个哑变量。对于 $N=2$，有 $\Phi\sim\diag(a,-a)$，因此除去量子修正，$u_2=-a^2$。类似地，对于一般的 $N$，除去量子修正，$u_k$ 是变量 $a_1,\ldots,a_N$ 的 $k$ 次初等对称多项式。因此，我们的任务是确定 $(u_2,\ldots,u_N)$ 与 $(a_1,\ldots,a_N)$ 之间包含量子修正的映射。

\subsection{纯 $\SU(N)$ 理论}\label{sec:pureSUN}\index{PureSUN@纯 $\SU(N)$ 理论|hyperpage}
\subsubsection{曲线}
不再赘述，直接引入 Seiberg-Witten 曲线。首先，纯 $\SU(N)$ 理论的 Seiberg-Witten 曲线为
\begin{equation}
\Sigma:\qquad\frac{\Lambda^N}z+\Lambda^Nz=x^N+u_2x^{N-2}+\cdots+u_N\label{puresuNcurve}
\end{equation}
微分仍为 $\lambda=xdz/z$。紫外曲线 $C$ 只是一个以 $z$ 为复坐标的球面。在紫外曲线上每个点 $z$ 处，上述方程都有 $N$ 个解。因此，$\Sigma$ 是 $C$ 的 $N$ 叶覆盖。

\begin{figure}[h]
\[
\inc{suN-wbosonX}
\]
\caption{$\SU(N)$ 理论的 W 玻色子。\label{fig:SUN-W}}
\end{figure}
检验这条曲线能否重现半经典行为。通过
\begin{equation}
x^N+u_2x^{N-2}+\cdots+u_N=\prod_i(x-\underline a_i).
\end{equation}
引入变量 $\underline a_i$。

将紫外曲线上的 $A$ 闭链取为单位圆 $|z|=1$。由于 Seiberg-Witten 曲线是 $N$ 叶覆盖，可以将这条曲线提升到每个叶片上，并将其称为闭链 $A_i$。假定所处区域满足 $|\underline a_i|\sim\LambdaRG$，与 $i$ 无关，且 $\LambdaRG\gg\Lambda$。那么，式 \eqref{puresuNcurve} 的近似解为
\begin{equation}
x_i=\underline a_i+O(\Lambda).
\end{equation}
将 $\lambda=xdz/z$ 本身视为叶片的坐标更加方便。于是有
\begin{equation}
\lambda_i=\underline a_i\frac{dz}z+O(\Lambda).
\end{equation}
这一情形画在图 \ref{fig:SUN-W} 中。很容易计算 $\lambda$ 沿闭链 $A_i$ 的积分：
\begin{equation}
a_i:=\frac1{2\pi i}\oint_{A_i}\lambda=\underline a_i+O(\Lambda).\label{suNpureA}
\end{equation}

现在，可以在第 $i$ 个叶片与第 $j$ 个叶片之间悬挂一张环形膜。这一对象的质量为
\begin{equation}
|\frac1{2\pi i}\oint_{A_i}\lambda-\frac1{2\pi i}\oint_{A_j}\lambda|
=|\frac1{2\pi i}\oint_A(\lambda_i-\lambda_j)|=|a_i-a_j|.
\end{equation}
这重现了 W 玻色子的质量。

为了看到单极子，需要理解 $N$ 叶覆盖 $\Sigma\to C$ 的分支结构。将组合 $y=\Lambda^N(z+1/z)$ 视为一个坐标较为方便。于是，可以认为方程 \eqref{puresuNcurve} 确定了多项式图像
\begin{equation}
y=P(x)=x^N+u_2x^{N-2}+\cdots+u_N
\end{equation}
与水平直线
\begin{equation}
y=\Lambda^N(z+\frac1z)
\end{equation}
的交点，如图 \ref{fig:suNbranching} 所示。当然，需要将图复化，不过读者应该能够理解其思想。
\begin{figure}[h]
\[
\inc{suNbr}
\]
\caption{纯 $\SU(N)$ 理论中有 $(N-1)$ 对分支点。\label{fig:suNbranching}}
\end{figure}

显然，在 $y=\Lambda(z+1/z)$ 的 $(N-1)$ 个取值处，$N$ 个叶片中的两个相遇；每个这样的取值都在 $z$ 球面上对应一对分支点 $z_i^\pm$，满足 $z_i^+z_i^-=1$。注意，第 $i$ 个叶片与第 $(i+1)$ 个叶片在这一对分支点处相遇。因此，可以在这一对分支点之间悬挂一张圆盘形膜，如图 \ref{fig:SUNmono} 所示。
\begin{figure}[h]
\[
\inc{suN-monopole}
\]
\caption{纯 $\SU(N)$ 理论的单极子。\label{fig:SUNmono}}
\end{figure}

在满足 $|\underline a_i|\sim|\LambdaRG|\gg|\Lambda|$ 的半经典区域中，有
\begin{equation}
|z_i^+|\sim\frac{\LambdaRG^N}{\Lambda^N},\qquad
|z_i^-|\sim\frac{\Lambda^N}{\LambdaRG^N}.
\end{equation}
将连接 $z_i^+$ 与 $z_i^-$ 的路径称为 $B_i$。于是
\begin{align}
M_\text{monopole}&=|\frac1{2\pi i}\int_{B_i}(\lambda_i-\lambda_{i+1})|\\
&\sim|(a_i-a_{i+1})\frac1{2\pi i}\int_{\Lambda^N/\LambdaRG^N}^{\LambdaRG^N/\Lambda^N}\frac{dz}z|\\
&\sim|(a_i-a_{i+1})\frac{2N}{2\pi i}\log\frac{\LambdaRG}\Lambda|.
\end{align}
作出对应
\begin{equation}
\tau(\LambdaRG)=\frac{2N}{2\pi i}\log\frac{\LambdaRG}\Lambda.
\end{equation}
后，便重现了单极子的质量。它正确地重现了纯 $\SU(N)$ 理论的单圈跑动。

\subsubsection{红外规范耦合矩阵}\index{PeriodMatrix@周期矩阵|hyperpage}\index{Genus@亏格|hyperpage}\index{GaugeCouplingMatrix@规范耦合矩阵|hyperpage}
检验我们的曲线是否满足低能 $\U(1)^{N-1}$ 理论的耦合矩阵正定这一条件。为此，需要更好地理解 Seiberg-Witten 曲线 $\Sigma$ 的几何。它是 $C$ 的一个 $N$ 叶覆盖，具有 $2N-2$ 个二阶分支点 $z_i^\pm$，以及两个 $N$ 阶分支点 $z=0,\infty$。于是，$\Sigma$ 的亏格 $g$ 由 Riemann-Hurwitz 公式\footnote{这并不难推导。假设对曲线 $C$ 作三角剖分，具有 $V$ 个顶点、$E$ 条边和 $F$ 个三角形，并使所有分支点都位于顶点处。有 $\chi(C)=V-E+F$。可以直接将边和三角形提升到 $\Sigma$：得到 $NE$ 条边和 $NF$ 个三角形。但是，顶点数少于 $NV$。设每个顶点 $p_i$ 处的分支度为 $\deg p_i$。那么，$\Sigma$ 的三角剖分具有 $NV-\sum(\deg p_i-1)$ 个顶点。最终得到 $\chi(S)=N\chi(C)-\sum_i(\deg p_i-1)$。} 确定：
\begin{equation}
\chi(\Sigma)=N\chi(C)-(2N-2)-2(N-1)
\end{equation}
% [Tutor Note: p.127 脚注末尾原文作 "$\chi(S)=N\chi(C)-\sum_i(\deg p_i-1)$"，其中 $S$ 应为 $\Sigma$：该脚注正是在推导曲线 $\Sigma$ 的亏格，紧接着的正文也作 $\chi(\Sigma)$。正文照原文排 $S$。]
其中 $\chi(\Sigma)=2-2g$ 和 $\chi(C)=2$ 是各曲面的欧拉示性数。得到 $g=N-1$。一维闭链的一组基由 $(2N-2)$ 条闭链 $A_i$ 和 $\tilde B^i$（$i=1,\ldots,N-1$）组成，其交数为
\begin{equation}
A_i\cdot A_j=0=\tilde B^i\cdot\tilde B^j,\qquad
A_i\cdot\tilde B^j=\delta_i^j.
\end{equation}
这里的点积计数两条一维闭链的交点数目，并包含符号。所得到的一组闭链画在图 \ref{fig:he} 中。

\begin{figure}[h]
\[
\inc{suNhyperellipticform}
\]
\caption{纯 $\SU(N)$ 理论的 Seiberg-Witten 曲线上的闭链 $A_i$ 与 $\tilde B_i$。\label{fig:he}}
\end{figure}

图 \ref{fig:SUNmono} 与图 \ref{fig:he} 的画法颇为不同。从 $A_1$ 到 $A_{N-1}$ 的闭链可以直接对应起来。就全纯形式的线积分而言，有
\begin{equation}
A_N=-A_1-A_2\cdots-A_{N-1}.
\end{equation}
相应地，由式 \eqref{suNpureA} 定义的变量 $a_i$ 并不线性独立，而且有
\begin{equation}
a_N=-a_1-\cdots-a_{N-1}.
\end{equation}

图 \ref{fig:SUNmono} 中的组合 $B_i-B_{i+1}$ 与 $A_i$ 正向相交，与 $A_{i+1}$ 负向相交。因此有
\begin{equation}
B_i-B_{i+1}=\tilde B^i-\tilde B^{i+1}.
\end{equation}
等价地说，$\tilde B^i$ 是以不依赖于 $i$ 的方式将开路径 $B_i$ 补全而得到的闭合一维闭链。随后在曲线上定义
\begin{equation}
a_D^i:=\frac1{2\pi i}\oint_{\tilde B^i}\lambda.
\end{equation}
考虑
\begin{equation}
\tau^{ij}:=\frac{\partial a_D^i}{\partial a_j}
=X_D^{ik}(X^{-1})^j_k\qquad
\end{equation}
其中
\begin{equation}
X_i^k:=\frac{\partial a_i}{\partial u_k},\qquad
X_D^{jk}:=\frac{\partial a_D^j}{\partial u_k}.
\end{equation}
定义
\begin{equation}
\omega_k=\frac\partial{\partial u_k}\lambda\Big|_{\text{$z$ 固定}},
\end{equation}
得到
\begin{equation}
\tau^{ij}=X_D^{ik}(X^{-1})^j_k\qquad\text{其中}\quad
X_i^k=\oint_{A_i}\omega_k,\quad
X_D^{jk}=\oint_{\tilde B_j}\omega_k.\label{periodmatrix}
\end{equation}
可以检验，$\omega_{2,3,\ldots,N}$ 构成 $\Sigma$ 上全纯非奇异一形式的一组基。以这种方式构造的矩阵 $\tau^{ij}$，在数学上称为 $\Sigma$ 的周期矩阵，并且已知满足
\begin{equation}
\tau^{ij}=\tau^{ji},\qquad\Im\tau^{ij}\ \text{正定}.\label{wzw}
\end{equation}
由第一个条件可知，局部存在函数 $F(a_i)$，使得
\begin{equation}
a_D^i=\frac{\partial F}{\partial a_i},\qquad
\tau^{ij}=\frac{\partial^2F}{\partial a_i\partial a_j}.
\end{equation}
这说明，将以这种方式定义的 $a_i$、$a_D^i$ 与 $\U(1)^{N-1}$ 规范理论低能描述中出现的 $a_i$ 对应起来是合理的。逆规范耦合矩阵由 $\Im\tau^{ij}$ 给出，其正定性由数学关系 \eqref{wzw} 保证。

\subsection{含有基本表示味的 $\SU(N)$ 理论}
\subsubsection{$N_f=1$}
接下来，考虑含有一个裸质量为 $\mu$ 的味 $(Q,\tilde Q)$ 的 $\SU(N)$ 理论。曲线为
\begin{equation}
\Sigma:\qquad\frac{\Lambda^{N-1}(x-\mu)}z+\Lambda^Nz=x^N+u_2x^{N-2}+\cdots+u_N.
\end{equation}
回顾半经典分析，当 $a_i\sim\mu$ 时会出现一个轻的带荷超多重态。我们来检验上述曲线能否重现这一行为。

首先，与前面一样引入 $\underline a_i$，并考虑所有 $|\underline a_i|$ 都远大于 $|\Lambda|$ 的半经典区域。紫外曲线上的 $A$ 闭链仍然取为 $|z|=1$。于是，与纯理论的情形一样，得到 $a_i\sim\underline{a_i}+O(\Lambda)$。

为了看到弱耦合区域中的额外奇点，定义 $\tilde z=z/\Lambda^{N-1}$。此时曲线为
\begin{equation}
\frac{x-\mu}{\tilde z}+\Lambda^{2N-1}\tilde z=x^N+u_2x^{N-2}+\cdots+u_N,
\end{equation}
在极弱耦合极限中，可以近似为
\begin{equation}
\frac{x-\mu}{\tilde z}=x^N+u_2x^{N-2}+\cdots+u_N=\prod(x-\underline a_i).
\end{equation}
当 $\underline a_i=\mu$ 时，方程发生因式分解，曲线分成两部分；否则，曲线是 $z$ 球面的一个光滑 $N$ 重覆盖。这表明，当 $\underline a_i=\mu$ 时，Seiberg-Witten 曲线上的一条一维闭链收缩，悬挂在那里的膜产生一个无质量超多重态，见图 \ref{fig:connection}。

还可以检验单圈跑动。大 $z$ 区域中的分支点 $z_i^+$ 不变，因为 $N_f=1$ 曲线在大 $z$ 区域的结构本身相对于纯理论曲线没有改变。因而
\begin{equation}
z_i^+\sim(\LambdaRG/\Lambda)^N.
\end{equation}
在小 $z$ 区域中，分支点位于 $\Lambda^{N-1}x/z$ 与 $P(x)$ 同阶的位置附近。假设 $|x|\sim|\underline a_i|\sim|\LambdaRG|$，便有
\begin{equation}
z_i^-\sim(\Lambda/\LambdaRG)^{N-1}.
\end{equation}
于是，单极子的质量为
\begin{align}
M_\text{monopole}&=|\frac1{2\pi i}\int_{B_i}\lambda|\\
&\sim|(a_i-a_{i+1})\frac1{2\pi i}\int_{\Lambda^{N-1}/\LambdaRG^{N-1}}^{\LambdaRG^N/\Lambda^N}\frac{dz}z|\\
&\sim|(a_i-a_{i+1})\frac{2N-1}{2\pi i}\log\frac{\LambdaRG}\Lambda|.
\end{align}
这给出
\begin{equation}
\tau(\LambdaRG)=\frac{2N-1}{2\pi i}\log\frac{\LambdaRG}\Lambda
\end{equation}
正如所期望的那样。

\subsubsection{一般味数}
更一般地，可以考虑如下曲线：
\begin{multline}
\Sigma:\qquad
\frac{\Lambda^{N-N_L}\prod_{i=1}^{N_L}(x-\mu_i)}z
+z\Lambda^{N-N_R}\prod_{i=1}^{N_R}(x-\mu'_i)\\
=x^N+u_2x^{N-2}+\cdots+u_{N-1}x+u_N\label{suNNf}
\end{multline}
其中 $N_L,N_R\le N$。当 $N_L=N_R=N$ 时，需要像含有四个味的 $\SU(2)$ 曲线 \eqref{naive_nf_4} 中那样引入复数 $f$、$f'$：
\begin{multline}
\Sigma:\qquad
f\cdot\frac{\prod_{i=1}^N(x-\mu_i)}z
+f'\cdot z\prod_{i=1}^N(x-\mu'_i)\\
=x^N+u_2x^{N-2}+\cdots+u_{N-1}x+u_N.\label{suN-saturated}
\end{multline}
还需要区分曲线中的质量参数与 BPS 质量公式中的质量参数；对于含有四个味的 $\SU(2)$ 理论，这一点已在第 \ref{sec:ident} 小节详细研究过。以下主要讨论味数少于 $2N-1$ 的情形。

考虑所有 $\mu_i$ 和 $\mu'_i$ 都很小的情形。进一步考虑满足 $|\underline a_i|\gg\Lambda$ 的区域。与往常一样，得到 $a_i=\underline a_i+O(\Lambda)$。分支点位于
\begin{equation}
|z_i^+|\sim\frac{\LambdaRG^{N-N_R}}{\Lambda^{N-N_R}},\qquad
|z_i^-|\sim\frac{\Lambda^{N-N_L}}{\LambdaRG^{N-N_L}}.
\end{equation}
于是有
\begin{align}
M_\text{monopole}&\sim|(a_i-a_{i+1})\frac1{2\pi i}\int_{\Lambda^{N-N_L}/\LambdaRG^{N-N_L}}^{\LambdaRG^{N-N_R}/\Lambda^{N-N_R}}\frac{dz}z|\\
&\sim|(a_i-a_{i+1})\frac{2N-(N_L+N_R)}{2\pi i}\log\frac{\LambdaRG}\Lambda|,
\end{align}
因此，单圈跑动为
\begin{equation}
\tau(\LambdaRG)=\frac{2N-(N_L+N_R)}{2\pi i}\log\frac{\LambdaRG}\Lambda.
\end{equation}

在另一区域 $|\mu_i|,|\underline a_i|\gg\Lambda$ 中，可以利用重定义的技巧寻找库仑支上的奇点。例如，定义 $\tilde z=z/\Lambda^{N-N_L}$，曲线变为
\begin{equation}
\frac{\prod_{i=1}^{N_L}(x-\mu_i)}{\tilde z}
+\tilde z\Lambda^{2N-N_R-N_L}\prod_{i=1}^{N_R}(x-\mu'_i)
=x^N+u_2x^{N-2}+\cdots+u_{N-1}x+u_N.
\end{equation}
此时可以取极限 $\Lambda\to0$，得到
\begin{equation}
\frac{\prod_{i=1}^{N_L}(x-\mu_i)}{\tilde z}=\prod_{i=1}^N(x-\underline a_i).
\end{equation}
这意味着，只要对某些 $i$ 和 $s=1,\ldots,N_L$ 有 $\underline a_i=\mu_s$，曲线就分成两部分，因为方程可以因式分解。对变量 $w=1/z$ 也可以作同样的处理。于是，当对某些 $i$ 和 $s=1,\ldots,N_R$ 有 $\underline a_i=\mu'_s$ 时，也会出现奇点。总之，这些结果重现了含有 $N_f=N_R+N_L$ 个基本表示超多重态的 $\SU(N)$ 理论的半经典弱耦合物理。图 \ref{fig:suNgeneral} 概括了这一情形。

\begin{figure}[h]
\[
\inc{suN-flavors}
\]
\caption{含有味的 $\SU(N)$ 理论。\label{fig:suNgeneral}}
\end{figure}
有一个用坐标 $z$ 描述的球面 $C$。曲线 $\Sigma$ 是 $C$ 的 $N$ 叶覆盖。一张 M5 膜包裹 $\Sigma$。我们将存在于 $C$ 上的六维理论称为 $\SU(N)$ 型 $\cN{=}(2,0)$ 理论。粗略地说，它来自 $N$ 张重合的 M5 膜。

将 $\Arg z$ 视为第六个方向 $x_6$，将 $\log|z|$ 视为第五个方向 $x_5$。沿 $x_6$ 方向约化，得到线段上的五维理论。这一五维理论是规范群为 $\SU(N)$ 的最大超对称杨--米尔斯理论。曲线中的项
\begin{equation}
\frac{\Lambda^{N-N_L}\prod_{i=1}^{N_L}(x-\mu_i)}z\label{Lb}
\end{equation}
可以视为定义了第五个方向左端的某种边界条件。我们认为，它在该处给出 $N_L$ 个 $\SU(N)$ 基本表示超多重态。类似地，项
\begin{equation}
z\Lambda^{N-N_R}\prod_{i=1}^{N_R}(x-\mu'_i)\label{Rb}
\end{equation}
被视为在该处具有 $N_R$ 个基本表示超多重态的边界条件。再沿第五个方向约化理论，就得到总共含有 $N_f=N_L+N_R$ 个基本表示超多重态的 $\SU(N)$ 规范理论。我们看到，边界条件的影响大致在
\begin{equation}
\log|z_R|\sim(N-N_R)\log\frac{\LambdaRG}\Lambda<0,\qquad
\log|z_L|\sim(N-N_L)\log\frac\Lambda{\LambdaRG}>0.
\end{equation}
时变得显著。在五维杨--米尔斯理论中，有单极子弦，其端点位于 $|z_R|$ 与 $|z_L|$ 附近。从四维观点看，$\log|z_L|/|z_R|$ 控制单极子的质量，进而给出理论的单圈跑动。

注意，从四维观点看，将 $N_f$ 分成 $N_R$ 和 $N_L$ 是颇为任意的。事实上，通过重定义 $z$，很容易将曲线化为如下形式：
\begin{equation}
\frac{\Lambda^{2N-N_f}\prod_i^{N_f}(x-\mu_i)}z
+z=x^N+u_2x^{N-2}+\cdots u_N
\end{equation}
其中定义了 $\mu_{N_L+i}:=\mu'_i$。在这一形式中，交换全部 $N_f$ 个质量参数的对称性是显然的。

然而，从高维观点看，有时保留 Seiberg-Witten 曲线 $\Sigma$ 的方程作为 $x$ 的多项式时次数为 $N$ 的形式，会更方便。这保证了 $\Sigma$ 始终是紫外曲线 $C$ 的 $N$ 叶覆盖。从次数上说，这一条件意味着，式 \eqref{Lb} 和式 \eqref{Rb} 这样的边界条件的次数也必须小于或等于 $N$。这就施加了约束 $N\ge N_{L,R}$，从而有 $2N\ge N_f$。这正是理论渐近自由或渐近共形的条件。
\subsection{$\SO(2N)$ 理论}\label{sec:SO}
现在简要讨论 $\SO(2N)$ 规范理论。

\subsubsection{半经典分析}
矢量多重态中的标量 $\Phi$ 是一个 $2N\times2N$ 反对称矩阵。将超多重态记为 $(Q_i^a,\tilde Q^i_a)$，其中 $a=1,\ldots,2N$，$i=1,\ldots,N_f$。考虑由
% [Tutor Note: p.132 原文 "an $2N\times2N$ antisymmetric matrix" 中冠词 "an" 应为 "a"（原文语法笔误）；另 p.125 原文 "the equation (11.2.1) can be thought of determining" 应为 "thought of as determining"（原文语法笔误）。两处均不影响中文语义，按正常语义译出。]
\begin{equation}
[\Phi,\Phi^\dagger]=0.
\end{equation}
给出的超对称真空分支。由于 $\Phi$ 反对称，对角化的结果为
\begin{equation}
\Phi=\diag(a_1,\ldots,a_N,-a_1,\ldots,-a_N).
\end{equation}
一般而言，规范群破缺到 $\U(1)^N$。规范不变组合由下式给出：
\begin{equation}
x^{2N}+u_2x^{2N-2}+u_4x^{2N-4}+\cdots+u_{2N}=\det(x+\Phi)
\end{equation}
其中 $x$ 是一个哑变量。注意，由于反对称性，奇次幂自动消失。事实上，由条件
\begin{equation}
u_{2N}=\tilde u_N{}^2,\qquad\tilde u_N=a_1a_2\ldots a_N
\end{equation}
定义的 $\tilde u_N$，在 $\SO(2N)$ 下也不变，但在 $\OO(2N)$ 下并不保持不变。

W 玻色子的质量为
\begin{equation}
|{}\pm a_i\pm a_j|
\end{equation}
其中 $i\neq j$。类似地，单极子的质量为
\begin{equation}
|\tau(\pm a_i\pm a_j)|.
\end{equation}

展开超势
\begin{equation}
\sum_i(Q_i\Phi\tilde Q^i+\mu_iQ_i\tilde Q^i),
\end{equation}
在经典层次上可知，当对某些 $i$ 和 $s$ 有 $\mu_s=\pm a_i$ 时，会出现一个在某个 $\U(1)$ 规范场下带荷的无质量超多重态。

单圈跑动为
\begin{equation}
\Lambda\frac d{d\Lambda}\tau=-\frac1{2\pi i}(2(2N-2)-2N_f).
\end{equation}
因此，当 $N_f<2N-2$ 时理论是渐近自由的，当 $N_f=2N-2$ 时是渐近共形的。

\subsubsection{纯 $\SO(2N)$ 理论}\label{sec:pureSON}\index{PureSO2N@纯 $\SO(2N)$ 理论|hyperpage}
纯理论的 Seiberg-Witten 曲线为
\begin{equation}
x^2(\frac{\Lambda^{2N-2}}z+\Lambda^{2N-2}z)
=x^{2N}+u_2x^{2N-2}+u_4x^{2N-4}+\cdots+u_{2N}\label{pureSOcurve}
\end{equation}
微分为 $\lambda=xdz/z$。它是紫外曲线 $C$ 的 $2N$ 叶覆盖，而 $C$ 只是一个以 $z$ 为复坐标的球面。求解这一方程，得到 $2N$ 个局部解 $\pm x_i(z)$。相应地，定义 $\pm\lambda_i=\pm x_i(z)dz/z$。

研究弱耦合区域。通过
\begin{equation}
x^{2N}+u_2x^{2N-2}+u_4x^{2N-4}+\cdots+u_{2N}
=\prod_{i=1}^N(x^2-\underline a_i{}^2).
\end{equation}
引入 $\underline a_i$。我们感兴趣的是满足 $|\underline a_i|\gg|\Lambda|$ 的区域。

\begin{figure}[h]
\[
\inc{so2N-wbosonX}
\]
\caption{$\SO(2N)$ 理论的 W 玻色子。\label{fig:soN}}
\end{figure}

在紫外曲线上 $|z|=1$ 处画出 $A$ 闭链，见图 \ref{fig:soN}。沿 $A$ 闭链，方程 \eqref{pureSOcurve} 的近似解为
\begin{equation}
x_i(z)=\underline a_i+O(\Lambda).
\end{equation}
将 $C$ 上的 $A$ 闭链提升到 $\Sigma$ 的各叶片上，得到 $N$ 对闭链 $\pm A_i$。于是
\begin{equation}
a_i=\frac1{2\pi i}\oint_{A_i}\lambda
=\frac1{2\pi i}\oint\lambda_i=\underline a_i+O(\Lambda).
\end{equation}
现在，可以在各叶片之间悬挂环形膜。它们显然具有质量
\begin{equation}
|{\pm a_i\pm a_j}|.
\end{equation}
我们发现，需要施加约束：M2 膜不能悬挂在第 $i$ 个叶片与第 $(-i)$ 个叶片之间，以禁止质量为 $|{\pm2a_i}|$ 的 W 玻色子。至于单极子，分支点大致位于
\begin{equation}
z^+\sim\left(\frac{\LambdaRG}\Lambda\right)^{2N-2},\quad
z^-\sim\left(\frac\Lambda{\LambdaRG}\right)^{2N-2}.
\end{equation}
于是，可以像 $\SU(N)$ 规范理论中那样近似计算单极子质量，得到
\begin{equation}
\sim|(a_i-a_j)\frac1{2\pi i}\int_{z^-}^{z^+}\frac{dz}z|
=|(a_i-a_j)\frac{2(2N-2)}{2\pi i}\log\frac{\LambdaRG}\Lambda|.
\end{equation}
由此可知，跑动耦合为
\begin{equation}
\tau(\LambdaRG)=\frac{2(2N-2)}{2\pi i}\log\frac{\LambdaRG}\Lambda,
\end{equation}
正确地重现了单圈分析。

研究低能耦合矩阵 $\tau^{ij}$。分支点位于 $z=0,\infty$，以及 $z$ 球面上一般位置的 $N$ 对点处。在 $z=\infty$ 处，有 $N-2$ 个解的行为为 $x\sim z^{1/(2N-2)}$，两个解的行为为 $x\sim z^{-1/2}$。因此，这里相当于一个分支度为 $2N-2$ 的分支点和另一个分支度为 $2$ 的分支点。$z=0$ 处的分支结构相同。接下来，考虑后一类的 $N$ 对分支点中的一对。当叶片 $i$ 与 $j$ 在那里相遇时，叶片 $-i$ 与 $-j$ 也同时相遇。将它们稍微分开，就得到总共 $4N$ 个分支度为 $2$ 的分支点。利用 Riemann-Hurwitz 定理，有
\begin{equation}
\chi(\Sigma)=2N\chi(C)-2(2N-3)-2-4N.
\end{equation}
因此，Seiberg-Witten 曲线的亏格为 $g=2N-1$。于是，$\Sigma$ 上独立的一维闭链可以标记为 $\tilde A_1,\ldots,\tilde A_{2N-1}$ 和 $\tilde B^1,\ldots,\tilde B^{2N-1}$，交数为
\begin{equation}
\tilde A_i\cdot\tilde A_j=0=\tilde B^i\cdot\tilde B^j,\qquad
\tilde A_i\cdot\tilde B^j=\delta_i^j.
\end{equation}
注意，曲线 $\Sigma$ 具有通过 $x\to-x$ 作用的 $\bZ_2$ 对称性。在这一对称性下，微分是奇的：$\lambda\to-\lambda$。相应地，只有在这一 $\bZ_2$ 作用下为奇的一维闭链 $L$，才可能有 $\oint_L\lambda\neq0$。将紫外曲线 $C$ 上的 $A$ 闭链提升到 $\Sigma$ 而得到的闭链 $A_i$（$i=1,\ldots,N$），确实是奇的。按式 \eqref{periodmatrix} 计算的周期矩阵 $\tau^{ij}$ 是一个 $(2N-1)\times(2N-1)$ 矩阵，它是对称的，且虚部正定。限制到 $\bZ_2$ 作用下为奇的子空间，就得到一个 $N\times N$ 矩阵，它同样是对称的，且虚部正定。

\subsubsection{含有矢量表示味的 $\SO(2N)$ 理论}
含有一个 $2N$ 维表示超多重态的 $\SO(2N)$ 理论，其曲线为
\begin{equation}
x^2(\frac{\Lambda^{2N-4}(x^2-\mu^2)}z+\Lambda^{2N-2}z)
=x^{2N}+u_2x^{2N-2}+u_4x^{2N-4}+\cdots+u_{2N}.
\end{equation}
我们只来检验，当对某个 $i$ 有 $a_i=\pm\mu$ 时，库仑支上会出现奇点。与往常一样，假设 $|\underline a_i|,|\mu|\gg|\Lambda|$，作重定义 $\tilde z=z/\Lambda^{2N-4}$，并对曲线取极限：
\begin{equation}
x^2\frac{(x^2-\mu^2)}z
=x^{2N}+u_2x^{2N-2}+u_4x^{2N-4}+\cdots+u_{2N}.
\end{equation}
当对某个 $i$ 有 $\pm\underline a_i=\mu$ 或 $\underline a_i=0$ 时，这个方程可以因式分解。后一种选择不满足假设 $|\underline a_i|\gg|\Lambda|$。因此，在弱耦合区域中，当 $\pm\underline a_i\sim\mu$ 时出现奇点。

一般而言，含有 $N_f=N_R+N_L$ 个矢量表示超多重态的 $\SO(2N)$ 理论，其曲线为
\begin{multline}
x^2(\frac{\Lambda^{2(N-N_R-1)}\prod_{i=1}^{N_R}(x^2-\mu_i^2)}z
+\Lambda^{2(N-N_L-1)}z\prod_{i=1}^{N_L}(x^2-\mu'_i{}^2))\\
=x^{2N}+u_2x^{2N-2}+u_4x^{2N-4}+\cdots+u_{2N}.\label{so2NNf}
\end{multline}
严格地说，这只适用于 $N_L+N_R<2N-2$。当 $N_L=N_R=N-1$ 时，需要用两个复数 $f$ 和 $f'$ 代替 $\Lambda$ 的幂，就像含有 $2N$ 个味的 $\SU(N)$ 理论的式 \eqref{suN-saturated} 中那样。

检验 $\mu_i=\mu'_i=0$ 时的单圈跑动。假定 $|\underline a_i|\gg|\Lambda|$。与往常一样，得到 $a_i=\underline a_i+O(\Lambda)$。紫外曲线上的分支点大致位于
\begin{equation}
z^+\sim\left(\frac{\LambdaRG}\Lambda\right)^{2N-2-2N_L},\quad
z^-\sim\left(\frac\Lambda{\LambdaRG}\right)^{2N-2-2N_R}.
\end{equation}
于是，可以像 $\SU(N)$ 规范理论中那样近似计算单极子质量，得到
\begin{equation}
\sim\left|(a_i-a_j)\frac1{2\pi i}\int_{z^-}^{z^+}\frac{dz}z\right|
=\left|(a_i-a_j)\frac{2(2N-2-2(N_L+N_R))}{2\pi i}\log\frac{\LambdaRG}\Lambda\right|.
\end{equation}
由此可知，跑动耦合为
\begin{equation}
\tau(\LambdaRG)=\frac{2(2N-2)-2(N_L+N_R)}{2\pi i}\log\frac{\LambdaRG}\Lambda,
\end{equation}
正确地重现了单圈分析。同样，理论渐近自由或共形的条件，与曲线方程左侧的次数小于或等于右侧的次数这一事实相联系。

\subsection{Argyres-Douglas 共形场论}\label{sec:ADAD}
研究本节所分析的各理论在库仑支上最奇异的点。

\subsubsection{纯 $\SU(N)$ 理论}
首先，取纯 $\SU(N)$ 理论的曲线：
\begin{equation}
\frac{\Lambda^N}z+\Lambda^Nz=x^N+\cdots+u_N
\end{equation}
微分为 $\lambda=xdz/z$。令 $z=1+\delta z$、$u_N=2\Lambda^N+\delta u_N$，并取 $\delta z$ 和 $\delta u_N$ 都很小的极限，得到
\begin{equation}
c\,\delta z^2=x^N+u_2x^{N-2}+\cdots+\delta u_N
\end{equation}
其中 $c$ 是一个不重要的常数。此时微分为 $\lambda=xd\delta z\sim\delta zdx$。引入 $\tilde z=1/x$，可以将这一极限中的曲线写为
\begin{equation}
c\,\lambda^2=\frac{1+u_2\tilde z^2+u_3\tilde z^3+\cdots+\tilde u_N\tilde z^N}{\tilde z^{N+4}}d\tilde z^2.
\end{equation}
注意，它与第 \ref{sec:generalAD} 小节中看到的曲线具有相同的形式；那里的曲线来自对
\begin{equation}
\lambda^2=\phi(\tilde z)
\end{equation}
的考虑，其中 $\phi(\tilde z)$ 是仅具有一个 $N+4$ 阶极点的二次微分，见图 \ref{fig:pureSUAD}。这与图 \ref{fig:XNYN} 中引入的理论 $Y_{N+4}$ 相同。因而有
\begin{equation}
AD_{N_f=0}(\SU(N))=Y_{N+4}.\label{@@@}
\end{equation}

\begin{figure}[h]
\[
\inc{N+IV}
\]
\caption{纯 $\SU(N)$ 理论最奇异的点。\label{fig:pureSUAD}}
\end{figure}

要求 $\lambda$ 具有标度维数 $1$，得到
\begin{equation}
[u_k]=\frac{2k}{N+2}.
\end{equation}
注意，有 $[u_k]+[u_{N+2-k}]=2$。此时回顾式 \eqref{generalCFTfeatures} 附近的讨论很有启发。考虑预势形变
\begin{equation}
\int d^4\theta\,u_ku_{N+2-k}
\end{equation}
其中 $d^4\theta$ 是手征 $\cN{=}2$ 超空间积分。当 $k\le N+2-k$ 时，有 $[u_k]\le1\le[u_{N+2-k}]$，因此将 $u_k$ 视为物理算符 $u_{N+2-k}$ 的形变参数。

取最简单的情形 $N=3$。该理论具有一个满足 $[u_3]=6/5$ 的算符，以及一个满足 $[u_2]=4/5$ 的相应参数。它们与来自含有一个味的 $\SU(2)$ 理论的 Argyres-Douglas 共形场论相同；事实上，曲线和微分也完全相同：
\begin{equation}
AD_{N_f=0}(\SU(3))=Y_7=AD_{N_f=1}(\SU(2)).
\end{equation}

\subsubsection{含有两个味的 $\SU(N)$ 理论}
接下来，考虑含有两个味的 $\SU(N)$ 理论。曲线为
\begin{equation}
(x-\mu_1)\frac{\Lambda^{N-1}}z+(x-\mu_2)\Lambda^{N-1}z=x^N+\cdots+u_N.
\end{equation}
已经研究过 $N=2$ 的情形，因此取 $N>3$。随后作展开
\begin{equation}
u_{N-1}=2\Lambda^{N-1}+\delta u_{N-1},\quad z=1+\delta z
\end{equation}
并取 $\delta z$、$\delta u_{N-1}$ 和 $\mu_{1,2}$ 都很小的极限。曲线为
\begin{equation}
c\,(x-\frac{\mu_1+\mu_2}2)\delta z^2+c'\,(\mu_1-\mu_2)\delta z
=x^N+\cdots+u_{N-2}x^2+\delta u_{N-1}x+u_N
\end{equation}
微分为 $\lambda=xd\delta z\sim\delta zdx$。这里，$c$ 和 $c'$ 是不重要的常数。

现在通过 $x=x'-(\mu_1+\mu_2)/2$ 定义 $x'$，将 $\delta z$ 平移为 $\delta z\to\delta z-(c'/c)(\mu_1-\mu_2)/(2x')$，并引入 $\tilde z=1/x'$。此时曲线为
\begin{equation}
\lambda^2=\frac{1+\tilde u_1\tilde z+\tilde u_2\tilde z^2+\cdots+\tilde u_N\tilde z^N+(\tilde\mu_1-\tilde\mu_2)^2\tilde z^{N+1}}{\tilde z^{N+3}}d\tilde z^2.
\end{equation}
这里，将各种不重要的数值常数吸收进带有波浪号的变量的定义中。

这就是曲线 $\lambda^2=\phi(z)$，其中 $\phi$ 具有一个 $N+3$ 阶极点和一个二阶极点，见图 \ref{fig:suNADwithF}。这就是图 \ref{fig:XNYN} 中引入的理论 $X_{N+3}$。有
\begin{equation}
AD_{N_f=2}(\SU(N))=X_{N+3}.\label{===}
\end{equation}
含有奇数个味的 $\SU(N)$ 理论最奇异的点，给出一个 $\cN{=}2$ 共形场论，文献 \cite{Cecotti:2013lda,Giacomelli:2013tia} 对其进行了分析。含有偶数个味、且 $N_f\ge4$ 的 $\SU(N)$ 理论最奇异的点，\emph{并不}给出一个 $\cN{=}2$ 共形场论；我们将在第 \ref{sec:applicationL} 小节看到这一点。

\begin{figure}[h]
\[
\inc{N+III-II}
\]
\caption{含有两个味的 $\SU(N)$ 理论最奇异的点。\label{fig:suNADwithF}}
\end{figure}

\subsubsection{纯 $\SO(2N)$ 理论}\index{PureSO2N@纯 $\SO(2N)$ 理论|hyperpage}
接下来，取纯 $\SO(2N)$ 理论
\begin{equation}
x^2(\frac{\Lambda^{2N-2}}z+\Lambda^{2N-2}z)
=x^{2N}+u_2x^{2N-2}+\cdots+u_{2N}.
\end{equation}
令
\begin{equation}
u_{2N-2}=2\Lambda^{2N-2}+\delta u_{2N-2},\qquad z=1+\delta z
\end{equation}
并取 $\delta u_{2N-2}$、$\delta z$ 都很小的极限。曲线为
\begin{equation}
c\,\delta z^2=x^{2N-2}+\cdots+\delta u_{2N-2}+\frac{u_{2N}}{x^2}
\end{equation}
其中 $c$ 是一个不重要的常数。微分为 $\lambda=xd\delta z\sim\delta zdz$。用 $\tilde z=1/x^2$ 表示，曲线为
\begin{equation}
c\,\lambda^2=\frac{1+\tilde u_2\tilde z+\cdots+\tilde u_{2N-2}\tilde z^{N-1}+u_{2N}\tilde z^N}{\tilde z^{N+2}}d\tilde z^2.
\end{equation}
这同样是曲线 $\lambda^2=\phi(z)$，其中 $\phi$ 具有一个相当高的 $N+2$ 阶极点，以及一个二阶极点，见图 \ref{fig:so2NAD}。因此有
\begin{equation}
AD_{N_f=0}(\SO(2N))=X_{N+2}.\label{qqq}
\end{equation}

\begin{figure}[h]
\[
\inc{N+II-II}
\]
\caption{$\SO(2N)$ 理论最奇异的点。\label{fig:so2NAD}}
\end{figure}

现在，$\SU(4)$ 和 $\SO(6)$ 具有相同的李代数。利用式 \eqref{@@@} 和式 \eqref{qqq}，得到
\begin{equation}
Y_8=AD_{N_f=0}(\SU(4))=AD_{N_f=0}(\SO(6))=X_5.
\end{equation}
利用式 \eqref{===}，可知它们也等价于 $AD_{N_f=2}(\SU(2))$。这一组等价关系解释了式 \eqref{accidents} 中看到的结果。

\subsubsection{Argyres-Douglas 共形场论与希格斯支}\index{WeakGauging@弱规范化|hyperpage}\index{HiggsBranch@希格斯支|hyperpage}
含有 $N_f=2$ 个味的 $\SU(2)$ 理论具有形如 $\bC^2/\bZ_2$ 的希格斯支，但纯 $\SU(4)$ 理论在紫外并没有它。我们刚刚声称
\begin{equation}
AD_{N_f=0}(\SU(4))=AD_{N_f=2}(\SU(2)).
\end{equation}
这两点如何相容呢？下面的讨论概括了文献 \cite{Argyres:2012fu} 的内容。

注意，极限中的 Argyres-Douglas 理论具有一个标度维数为 $4/3$ 的算符 $u$、一个标度维数为 $2/3$ 的相应参数 $m$，以及一个标度维数为 $1$ 的额外质量参数 $\mu_1-\mu_2$。当将它实现为含有 $N_f=2$ 个味的 $\SU(2)$ 理论的极限时，显然，低能理论只有一个 $\U(1)$ 多重态，而 $\mu_1-\mu_2$ 是一个外部参数。

然而，当将同一理论实现为纯 $\SU(4)$ 理论的极限时，原来的低能理论具有 $\U(1)^3$ 矢量多重态，以及三个库仑支参数 $u_2$、$u_3$、$u_4$。已经看到，$\delta u_2$ 的标度维数为 $2/3$，$\delta u_4$ 的标度维数为 $4/3$，$\delta u_3$ 的标度维数为 $1$。因此，在这种实现中，极限 Argyres-Douglas 理论的质量参数 $\mu_1-\mu_2$ 被提升为一个 $\U(1)$ 多重态的真空期望值 $\delta u_3$。等价地说，极限理论的 $\SU(2)$ 味对称性中的 $\U(1)$ 子群，被动力学地弱规范化了，从而去除了希格斯支。

类似地，这里看到纯 $\SO(8)$ 理论与含有 $N_f=2$ 个味的 $\SU(3)$ 理论，都会给出共形场论 $X_6$。在第 \ref{sec:generalAD} 小节还知道，含有 $N_f=3$ 个味的 $\SU(2)$ 理论，在库仑支上也有一点，其低能极限由同一理论 $X_6$ 描述，见式 \eqref{accidents}。关于它们的希格斯支，也可以作与上面类似的研究。极限理论本身具有一个标度维数为 $3/2$ 的算符 $u$、一个维数为 $1/2$ 的相应形变参数 $m$，以及 $\SU(3)$ 味对称性的两个质量参数 $\mu_1-\mu_2$ 和 $\mu_1-\mu_3$。将其描述为含有 $N_f=3$ 个味的 $\SU(2)$ 理论的库仑支上的一点时，这些性质最为清楚。

在含有 $N_f=2$ 个味的 $\SU(3)$ 理论的描述中，有两个库仑支算符 $u_2$、$u_3$，味对称性 $\U(1)$ 部分的质量参数 $m$，以及 $\SU(2)$ 部分的质量参数 $\mu_1-\mu_2$。可以看到，$u_2$ 与 $u_3$ 的标度维数分别为 $1$ 和 $3/2$，$m$ 的标度维数为 $1/2$，$\mu_1-\mu_2$ 的维数为 $1$。因此，味对称性 $\SU(3)$ 的 $\U(1)$ 子群被弱规范化了。承担这一弱规范化的 $\U(1)$ 矢量多重态的真空期望值为 $u_2$。

在纯 $\SO(8)$ 理论的描述中，有四个库仑支算符 $u_2$、$u_4$、$u_6$、$u_8$，不过如前所述，$u_8=\tilde u_4^2$。在 Argyres-Douglas 点附近，$u_2$、$u_4$、$u_3$ 和 $\tilde u_4$ 的标度维数分别为 $1/2$、$1$、$3/2$ 和 $1$。可以看到，味对称性 $\SU(3)$ 的 $\U(1)^2$ 子群，由标量分量为 $u_4$ 和 $\tilde u_4$ 的两个 $\U(1)$ 矢量多重态弱规范化。$\SO(8)$ 的外自同构 $S_3$ 对维数为 $1$ 的算符 $u_4$ 与 $\tilde u_4$ 的作用，由宇称操作 $\tilde u_4\to-\tilde u_4$ 和作用于 $u_4$--$\tilde u_4$ 平面的 $120^\circ$ 旋转生成。这恰好就是味对称性 $\SU(3)$ 的 Weyl 群对满足 $\mu_1+\mu_2+\mu_3=0$ 的两个质量参数 $\mu_1$、$\mu_2$、$\mu_3$ 的作用方式。因此，$\SO(8)$ 的外自同构对称性可以与 $\SU(3)$ 味对称性的 Weyl 群对应起来。

\subsection{各种其他单规范群的 Seiberg-Witten 解}\label{sec:others}\index{HyperellipticCurve@超椭圆曲线|hyperpage}\index{SeibergWittenGeometry@Seiberg-Witten 几何|hyperpage}
最后，非常简要地提及各种其他规范理论的 Seiberg-Witten 解，以结束本节。首先，重写含有 $N_f$ 个味的 $\SU(N)$ 与 $\SO(2N)$ 理论的解。含有 $N_f$ 个味的 $\SU(N)$ 理论的 Seiberg-Witten 曲线为
\begin{equation}
\frac{\Lambda^N}z+z\Lambda^{N-N_f}\prod_{i=1}^{N_f}(x-\mu_i)
=x^N+u_2x^{N-2}+\cdots+u_{N-1}x+u_N
\end{equation}
而含有 $N_f$ 个味的 $\SO(2N)$ 理论的 Seiberg-Witten 曲线为
\begin{equation}
x^2\left[\frac{\Lambda^{2N-2}}z+\Lambda^{2N-2-2N_f}z\prod_{i=1}^{N_f}(x^2-\mu_i{}^2)\right]
=x^{2N}+u_2x^{2N-2}+u_4x^{2N-4}+\cdots+u_{2N}.
\end{equation}
这里，为了简单起见，省略了左侧乘在 $1/z$ 上的味项。完整表达式分别见式 \eqref{suNNf} 和式 \eqref{so2NNf}。

这里不作推导，直接给出其他经典规范群的 Seiberg-Witten 曲线。含有 $N_f$ 个味的 $\SO(2N+1)$ 理论的 Seiberg-Witten 曲线为
\begin{equation}
x\left[\frac{\Lambda^{2N-1}}{z^{1/2}}+\Lambda^{2N-1-2N_f}z^{1/2}\prod_{i=1}^{N_f}(x^2-\mu_i{}^2)\right]
=x^{2N}+u_2x^{2N-2}+u_4x^{2N-4}+\cdots+u_{2N},
\end{equation}
而含有 $N_f$ 个味的 $\Sp(N)$ 理论的 Seiberg-Witten 曲线为
\begin{multline}
\frac{\Lambda^{2N+2}}{z^{1/2}}+2c+\Lambda^{2N+2-2N_f}z^{1/2}\prod_{i=1}^{N_f}(x^2-\mu_i{}^2)\\
=x^2(x^{2N}+u_2x^{2N-2}+u_4x^{2N-4}+\cdots+u_{2N})
\end{multline}
其中 $c^2=\Lambda^{4N+4-2N_f}\prod_{i=1}^{N_f}(-\mu_i^2)$。微分始终只是 $\lambda=xdz/z$。同样省略了左侧乘在 $1/z$ 上的味项。

迄今为止的曲线始终可以写成
\begin{equation}
\frac{F(x)}z+z\tilde F(x)=P(x)\label{FP}
\end{equation}
其中 $F(x)$、$\tilde F(x)$ 和 $P(x)$ 是某些多项式。较早的文献更常将曲线和微分写成如下形式：
\begin{equation}
y^2=P(x)^2-4F(x)\tilde F(x),\qquad
\lambda=\frac x2d\log\frac{P(x)-y}{P(x)+y}.\label{yFP}
\end{equation}
为了联系式 \eqref{FP} 与式 \eqref{yFP}，注意方程 \eqref{yFP} 意味着组合
\begin{equation}
\zeta_\pm=\frac12(P(x)\pm y)
\end{equation}
满足
\begin{equation}
\zeta_++\zeta_-=P(x),\qquad\zeta_+\zeta_-=F(x)\tilde F(x).
\end{equation}
与式 \eqref{FP} 比较，得到
\begin{equation}
\zeta_+=\frac{F(x)}z,\qquad\zeta_-=z\tilde F(x).
\end{equation}
这也解释了式 \eqref{yFP} 中给出的微分。

在某些情形下，这种较早的形式在数学上更容易处理，因为函数 $y(x)$ 的分支割线至多为二阶。在数学上，这样的黎曼曲面称为超椭圆曲面，与更一般的黎曼曲面相比，具有一些特殊性质。尽管如此，本讲义其余部分采用的形式具有更明确的物理意义，通常也更有用。

关于上述经典规范群的曲线，可以在例如文献 \cite{Nekrasov:2004vw} 中找到很好的总结。对于例外规范群，情况更为复杂。虽然可以写出 Seiberg-Witten 曲线，但更自然的是研究 Seiberg-Witten 几何；它是一个以紫外曲线 $C$ 为底的复三维纤维化空间。对于纯理论，文献 \cite{Hollowood:1997pp} 给出了非常好的论述。对于含有物质超多重态的情形，较新的参考文献是 \cite{Tachikawa:2011yr}。
